\label{sec:5}

% Conclusions provide a clear and thorough summary, with interesting future work suggested.
% Thoughtful personal reflection on the lessons learned.

This dissertation investigated whether the AI Mother Tongue architecture can support emergent discrete communication in a complex multi-agent POMDP when combined with a high-throughput MAPPO training pipeline and introduced a simpler alternative - WhisperWise. 
The results showed that whilst the AI Mother Tongue architecture establishes meaningful communication in simpler environments, it struggles to scale to more demanding tasks.
In contrast, WhisperWise successfully scales to these complex environments, producing communication.

\section{Work Completed}
\label{sec:5.1}

First, I implemented a configurable vectorised, high-throughput MARL environment capable of achieving over 100K steps per second.
The simulation modelled a discrete two-dimensional environment populated by agents and stations.
A lightweight real-time visualisation was created to enable inspection of the rollout phase.

Second, a MAPPO-based MARL training script under a centralised training, decentralised execution framework was implemented.
It uses a shared actor, a global critic, an LSTM module, a custom rollout buffer for efficient storage and GAE calculations, and a communication protocol interface that supports hot-swapping.
To reduce training time, a separate script was introduced to pre-train the actor-critic networks using imitation learning.

Third, I faithfully followed and extended the AI Mother Tongue architecture \cite{AIMotherTongue}.
Pre-trained VQ-VAE and SQ-VAE (and an exploratory HQ-VAE variant, detailed in Appendix \ref{sec:hq-vae}) were trained on optimal rollout phases, injected into the actor network, and frozen before MARL training.
To establish rigorous baselines and test novel improvements, alternative protocols were also developed, including Discrete Communication, No Communication and WhisperWise.

Fourth, a quantitative evaluation of communication protocols was performed.
The analysis compared task reward across the communication protocols.
This directly answered the research question proposed at the start of this dissertation.

Finally, the semantics of communication and memory were analysed.
By collecting and visualising conditional distributions, such as the probability of a partner's target given a code and the probability of a partner's next action given a code, the dissertation examined whether the message correlates with goals or with the next action.
This analysis showed the agent's communication carried interpretable meaning.

Ultimately, the evaluation proved that the AI Mother Tongue architecture successfully enables robust, semantically meaningful discrete communication, outperforming both on-policy discrete and no-communication baselines by overcoming the Communication Vacuum Equilibrium in simple tasks.
However, the AI Mother Tongue architecture overcomplicates the solution to the "Joint Exploration Problem".
The reflection strategies teach agents to use the communication channel effectively.
The frozen encoder encrypts the external observations, stunting the agents from learning the task reward.
Consequently, this dissertation presents WhisperWise: a new, streamlined protocol that merges on-policy discrete communication with reflection strategies.
For basic tasks, WhisperWise significantly outperforms the AI Mother Tongue architecture (mean difference of +0.0413, $p_{adj} < 0.001$; see Section \ref{sec:4.8}) and achieves a steeper, more efficient semantic compression.
WhisperWise can also robustly scale to complex 15x15 environments where AI Mother Tongue completely fails.
This proves that reflection strategies, not the semantic bottlenecks, are the primary component of successful emergent communication.

\section{Lessons Learnt}
\label{sec:5.2}

I have learned the importance of logging and checkpointing; when running a program that takes over seventeen hours to complete, waiting until the end to find it overfitted hours ago is catastrophic without these systems in place.
In terms of MARL communication, I learned that the complexity of an emergent language is fundamentally bounded by the environment that necessitates it.
Agents will reliably find the most compressed, 'lazy' communication strategy that satisfies the reward function.
% Maybe another lesson


\section{Social and Ethical Considerations}
\label{sec:5.3}

Integrating these trained models into real-world settings carries social and ethical implications.
A significant concern is the "black box" nature of MARL communication, creating a trust deficit in applications.
The analysis tools produced take a step further towards AI transparency, reducing this trust deficit.

In this project, data is generated from a built-in simulation, but for integration into real-world settings, real data must be used.
For example, the expert agent would likely be human, triggering the application of strict data protection laws (such as GDPR \cite{gdpr}).


\section{Future Work}
\label{sec:5.4}

This dissertation suggests several promising future directions.

For communication, each agent's code is broadcast to all other agents.
For simulations with more than two agents, this will inevitably lead to agents listening to code that is not relevant to their policy.
An improvement would be to apply an attention mechanism here to determine which codes are sent to which agents, similar to TarMAC \cite{tarmac}.
This attention mechanism has another key potential use.
When an agent observes the environment, they only see the closest visible entity.
This could be improved by using an attention mechanism to select which visible entity to observe.

A second direction for improvement is to evaluate whether the established communication protocols can scale to highly complex, hierarchical environments (such as 50x50 grids with agents passing ingredients to one another).
This would, in theory, push the HQ-VAE to establish macro- and micro-strategies (Appendix \ref{sec:hq-vae}).

Finally, a third direction is to integrate WhisperWise into a real-world application.

In conclusion, this work shows that the AI Mother Tongue architecture overcomplicates the solution to the "Joint Exploration Problem".
The reflection strategies teach agents to use the communication channel effectively, but the semantic bottleneck impedes their learning.
This dissertation introduces WhisperWise, a scalable and elegant solution that takes a step closer to developing artificial agents that can coordinate effectively and transparently.